\Askhsh{14477}{4}
\begin{ylh}{2}
    \item 6.2. Γραφική Παράσταση Συνάρτησης
    \item 6.3. Η Συνάρτηση $f(x)=ax+\beta$
\end{ylh}

Στο παρακάτω σύστημα συντεταγμένων η ευθεία $\varepsilon_1$ με $a_1 > 0$ και $\beta_1 > 0$ παριστάνει τη γραφική παράσταση της συνάρτησης $f(x) = a_1 x + \beta_1$ των ετήσιων δαπανών μιας εταιρείας, σε χιλιάδες ευρώ, στα $x$ χρόνια της λειτουργίας της.

Η ευθεία $\varepsilon_2$ με $a_2 > 0$ και $\beta_2 > 0$ παριστάνει τη γραφική παράσταση της συνάρτησης των ετήσιων εσόδων $g(x) = a_2 x + \beta_2$ της εταιρείας, σε χιλιάδες ευρώ, στα $x$ χρόνια της λειτουργίας της. Οι γραφικές παραστάσεις αναφέρονται στα δέκα πρώτα χρόνια λειτουργίας της εταιρείας.

\begin{center}
\begin{tikzpicture}
\begin{axis}[
    axis lines = middle,
    xlabel = {$x$ (χρόνια)},
    ylabel = {χιλ. \euro},
    xmin = 0, xmax = 11,
    ymin = 0, ymax = 80,
    xtick = {1,2,3,4,5,6,7,8,9,10},
    ytick = {10,20,30,40,50,60,70},
    grid = both,
    grid style = {gray!30},
    width = 10cm, height = 8cm,
    tick label style = {font=\small},
    every axis x label/.style = {at={(current axis.right of origin)}, anchor=south},
    every axis y label/.style = {at={(current axis.above origin)}, anchor=south},
]
    % ε1: δαπάνες f(x) = 4x + 10
    \addplot[very thick, secondarycolor, domain=0:10] {4*x + 10};
    \node[secondarycolor, above left] at (axis cs:8,12) {$\varepsilon_1$ (δαπάνες)};

    % ε2: έσοδα g(x) = 6x + 5
    \addplot[very thick, maincolor, domain=0:10] {6*x + 5};
    \node[maincolor, below] at (axis cs:5,59) {$\varepsilon_2$ (έσοδα)};

    % Σημείο τομής (2.5, 20)
    \addplot[only marks, mark=*, mark size=1.5pt, black] coordinates {(2.5,20)};
    \node[above left] at (axis cs:2.5,20) {$K$};
\end{axis}
\end{tikzpicture}\captionof{figure}{Άσκηση 14477}
\end{center}

\begin{alist}
\item Με τη βοήθεια των γραφικών παραστάσεων να εκτιμήσετε τα έσοδα και τις δαπάνες τον πέμπτο χρόνο λειτουργίας της εταιρείας.

\monades{4}
\item
\begin{rlist}
\item Να προσδιορίσετε τους τύπους των συναρτήσεων $f$, $g$ και να ελέγξετε αν οι εκτιμήσεις σας στο α) ερώτημα ήταν σωστές.

\monades{15}
\item Να βρείτε τις συντεταγμένες του σημείου τομής των ευθειών $\varepsilon_1$ και $\varepsilon_2$ και να τις ερμηνεύσετε στο πλαίσιο του προβλήματος.

\monades{6}
\end{rlist}
\end{alist}

